\section{Augmented Reality}

\begin{ejercicio}
Compare model-free and model-based tracking in AR. How do both approaches work, what are the differences, and for which real-world application would each approach be suitable?
\end{ejercicio}

\begin{ejercicio}
    Explain the concepts of ``markers'' and ``SLAM'' in AR. How do they differ and which technique (model-based or model-free) do each of them use?
\end{ejercicio}

\begin{ejercicio}
Name one of the key differences between Video See-Through (VST) and Optical See-Through (OST). Explain two advantages and two disadvantages of each. Name one scenario for which each would fit best.
\end{ejercicio}

\begin{ejercicio}
Explain the concepts of tracking and registration in Augmented Reality, detailing their role in overlaying virtual content onto the real world. Differentiate between geometric and photometric registration, providing a brief example for each. Why is low latency crucial for a good AR experience?
\end{ejercicio}

\begin{ejercicio}
List and explain two input modalities that can be useful to interact in AR. Name one advantage and one disadvantage for each.
\end{ejercicio}

\begin{ejercicio}
Which input modality (Rigid Object Tracing, Gestures, Pointing and Drawing, Touch) would you implement for this scenario? Name two reasons.
\end{ejercicio}

\begin{ejercicio}
In the lecture, various interaction techniques for augmented reality were discussed, such as gesture-based interaction, voice commands, and gaze monitoring. Consider the following use case and explain which interaction method you would recommend and why:
\begin{itemize}
    \item \ul{Use Case:} You are designing an AR application for elderly users that helps them perform light exercises at home. The app needs to be intuitive, minimize physical strain, and avoid requiring fine motor skills.
\end{itemize}
\end{ejercicio}

\begin{ejercicio}
Define what situated visualization in AR means, contrast it with conventional visualization methods, and provide two different techniques used for situated visualization along with an example use case for each.
\end{ejercicio}


\begin{ejercicio}
    Given these three scenarios and these three gestures, which gesture would you use for each scenario and why?
    \begin{itemize}
        \item Possible gestures: \textit{Pen gestures, Device gestures, Eye gestures, Hand gestures}
    \end{itemize}
\end{ejercicio}

\begin{ejercicio}
    During the lecture, gestural interaction techniques for AR were discussed. Explain two advantages and two disadvantages of using gestures as an input modality in AR.
\end{ejercicio}

\begin{ejercicio}
    Describe the concept of ``Reality-Virtuality Continuum''.
\end{ejercicio}

\begin{ejercicio}
    Explain the ``Midas Touch'' problem in AR. How can it be mitigated? Provide one example of a real-world application where this problem could arise and how it could be addressed.
\end{ejercicio}